% !TeX root = ../../main.tex
% sections/app/appendix_risk_sharing.tex

\chapter{完備市場における消費の共通化と貯蓄の個別化}
\label{chap:appendix_risk_sharing}

本付録では、国内完備市場とカルボ型価格設定を仮定したモデルにおいて、なぜ全ての家計の消費が共通化される一方、貯蓄（資産ポートフォリオ）は家計ごとに個別化されるのかを厳密に証明する。

\section*{定理9：リスク共有の結果}
\label{sec:appendix_risk_sharing_theorem9}

国内金融市場が完備である経済において、以下の関係が成立する。
\begin{enumerate}
    \item\label{itm:appendix_risk_sharing_theorem9_lambda}全ての国内家計\(h\in H\)の所得の限界効用は、いかなる時点\(t\)、いかなる状態\(j\)においても常に一致する。
    \begin{equation}
    \lambda_t^h=\lambda_t^{h'}\quad\forall h,h'\in H
    \label{eq:appendix_risk_sharing_theorem9_lambda_equality}
    \end{equation}
    \item\label{itm:appendix_risk_sharing_theorem9_consumption}全ての国内家計\(h\in H\)の総消費指数は、常に一致する。
    \begin{equation}
    c_t^{h\to W}=c_t^{h'\to W}\quad\forall h,h'\in H
    \label{eq:appendix_risk_sharing_theorem9_consumption_equality}
    \end{equation}
    \item\label{itm:appendix_risk_sharing_theorem9_portfolio}各家計が購入する次期のための資産ポートフォリオの総価値は、家計が当期に価格改定の機会を得たか否かによって異なるため、一般に一致しない。
\end{enumerate}

\section*{証明}
\label{sec:appendix_risk_sharing_proof}

\subsection*{1.所得の限界効用\(\lambda_t\)の一致（定理9-1の証明）}
\label{sec:appendix_risk_sharing_sub1}

\paragraph{ステップA：限界効用の比率の不変性}
\label{sec:appendix_risk_sharing_sub1_step_a}
家計の最適化行動は全ての時点・状態で成立するため、任意の2つの国内家計\(h\)と\(h'\)の国内コンティンジェント債券に関する一階の条件（FOC）は、常に成立する。
\begin{align}
\lambda_t^hq_{t+1}&=\beta_t^H\pi\lambda_{t+1}^h \label{eq:appendix_risk_sharing_sub1_foc_h} \\
\lambda_t^{h'}q_{t+1}&=\beta_t^H\pi\lambda_{t+1}^{h'} \label{eq:appendix_risk_sharing_sub1_foc_h_prime}
\end{align}
これら2式の比を取ると共通項が消去され、以下の関係が得られる。
\begin{equation}
\frac{\lambda_t^h}{\lambda_t^{h'}}=\frac{\lambda_{t+1}^h}{\lambda_{t+1}^{h'}}=k
    \label{eq:appendix_risk_sharing_sub1_constant_ratio}
\end{equation}
この比率\(k\)は、時間や状態に依存しない普遍的な定数である。

\paragraph{ステップB：定数\(k\)の値の特定}
\label{sec:appendix_risk_sharing_sub1_step_b}
第\ref{sec:model_overview}節で定義した「事的事実対称性」の仮定を用いる。初期時点\(s\)において、全ての家計は同一の選好をもち、かつ初期貯蓄がゼロ（\(d_s^h=d_s^{h'}=0\)）である。このとき、各家計が直面する最適化問題は数学的に完全に同一であるため、初期の所得の限界効用は全ての家計で一致しなければならない。
\begin{equation}
\lambda_s^h=\lambda_s^{h'}\quad\Longrightarrow\quad k=\frac{\lambda_s^h}{\lambda_s^{h'}}=1
\label{eq:appendix_risk_sharing_sub1_k_identity}
\end{equation}

\paragraph{ステップC：結論}
\label{sec:appendix_risk_sharing_sub1_step_c}
普遍的な定数が\(k=1\)であることから、将来のすべての時点・すべての状態において\(\lambda_t^h=\lambda_t^{h'}\)が成立する。

\subsection*{2.総消費指数\(c_t^{h\to W}\)の一致（定理9-2の証明）}
\label{sec:appendix_risk_sharing_sub2}
全ての家計の\(\lambda_t\)が一致するという結果を、消費に関するFOC（式\ref{eq:model_foc_consumption_home}）に適用する。
\begin{equation}
\lambda_t=\frac{1}{p_t^{H\to W}c_t^{h\to W}}
\label{eq:appendix_risk_sharing_sub2_lambda_def}
\end{equation}
左辺の\(\lambda_t\)と右辺の物価指数\(p_t^{H\to W}\)は全ての家計で共通であるため、総消費指数\(c_t^{h\to W}\)もまた、全ての家計間で完全に一致しなければならない。

\subsection*{3.ポートフォリオ購入総額の個別化（定理9-3の証明）}
\label{sec:appendix_risk_sharing_sub3}
家計\(h\)の予算制約式を、次期のために購入する資産ポートフォリオの総価値\(v_{t+1}^h\equiv\sum_{j'}q_{t+1}d_{t+1}^{h\to H}+e_tb_{t+1}^{h\to F*}\)について整理する。



\begin{equation}
v_{t+1}^h=\underbrace{\left(d_t^{h\to H}+(1+i_{t-1}^F)e_tb_t^{h\to F*}+(1-\tau_t^H)p_t^hy_t^h+t_t^H\right)}_{\text{当期の総収入}}-\underbrace{p_t^{H\to W}c_t^{h\to W}}_{\text{当期の消費支出}}
\label{eq:appendix_risk_sharing_sub3_portfolio_dynamics}
\end{equation}
定理9-2より消費支出は共通であるが、当期の総収入、特に労働所得\((1-\tau_t^H)p_t^hy_t^h\)はカルボ型の価格設定（\(p_t^h\)の異質性）によって家計ごとに異なる。支出が共通で収入が異なる以上、その差額を埋める資産保有額\(v_{t+1}^h\)は、各家計が直面したショックに応じて個別化される。（証明終）